\documentclass[11pt,a4paper]{article}

\usepackage[T1]{fontenc}
\usepackage[utf8]{inputenc}
\usepackage[french]{babel}
\usepackage[a4paper,margin=24mm]{geometry}
\usepackage{xcolor}
\usepackage{hyperref}
\usepackage{tabularx}
\usepackage{enumitem}

\definecolor{damerginavy}{HTML}{171D29}
\definecolor{damergigold}{HTML}{9C7540}
\definecolor{damergicream}{HTML}{F8F6F1}

\hypersetup{
  colorlinks=true,
  linkcolor=damerginavy,
  urlcolor=damergigold,
  pdftitle={DAMERGI.COM -- Documentation v0.7.1},
  pdfauthor={Ahmed-Farouk DAMERGI}
}

\setlist[itemize]{leftmargin=1.5em,itemsep=0.25em,topsep=0.4em}
\setlist[enumerate]{leftmargin=1.7em,itemsep=0.25em,topsep=0.4em}

\newcommand{\tech}[1]{\texttt{#1}}
\newcommand{\version}{0.7.1}

\begin{document}

\begin{center}
  {\color{damergigold}\rule{0.72\textwidth}{0.8pt}}\\[1.2em]
  {\Large\bfseries AFD -- AHMED-FAROUK DAMERGI}\\[1.1em]
  {\Huge\bfseries DAMERGI.COM}\\[0.5em]
  {\Large Portfolio Data \& Business Analyst}\\[1.2em]
  {\large Documentation technique et produit -- version \version}\\[0.8em]
  {\small Patch documentaire cumulatif de la version 0.7.0}\\[1.2em]
  {\color{damergigold}\rule{0.72\textwidth}{0.8pt}}
\end{center}

\vspace{1em}

\begin{quote}
\textbf{Statut documentaire.} Aucun source LaTeX antérieur n'était présent dans le dossier
\tech{docs/} au moment de cette évolution ; seule la documentation complète
\tech{damergi\_documentation\_v0.7.0.pdf} était disponible. Ce document conserve ce PDF
historique sans le modifier et constitue le patch versionné 0.7.1. Toutes les décisions,
contraintes et sections de la version 0.7.0 qui ne sont pas remplacées ci-dessous restent
applicables. En cas de reprise du projet, lire d'abord le PDF 0.7.0 puis ce source 0.7.1.
\end{quote}

\tableofcontents

\section{Objet de la version 0.7.1}

La version 0.7.1 enrichit l'écosystème animé de la section \emph{Expertise / Ce que je fais}
avec une septième branche consacrée au développement et au prototypage autour de la donnée.
Cette évolution montre la capacité à construire des interfaces, des outils internes, des
prototypes et des applications d'appui sans modifier le positionnement principal du portfolio :
\textbf{Data \& Business Analyst}.

La version ne crée aucune nouvelle expérience professionnelle et ne modifie ni les métriques
vérifiées, ni la timeline, ni ses modales, ni les comportements stables du hero. TypeScript,
React et Next.js sont notamment présentés comme des technologies démontrées par
DAMERGI.COM lui-même.

\section{Architecture conservée}

L'architecture décrite en version 0.7.0 reste inchangée :

\begin{verbatim}
app/page.tsx
  +-- Navbar / LoadingScreen / PageProgress
  +-- StoryHero
  +-- ExpertiseSection

data/
  +-- timeline.ts
  +-- timelineDetails.ts

app/globals.css
  +-- design system clair/sombre
  +-- responsive et reduced motion
\end{verbatim}

La nouvelle expertise est déclarée dans le contenu bilingue de
\tech{components/Expertise/ExpertiseSection.tsx}. Les chemins SVG et positions de nœuds
restent centralisés par clé d'expertise dans le même composant. Aucun changement n'est
nécessaire dans \tech{app/page.tsx}, \tech{data/timeline.ts},
\tech{data/timelineDetails.ts} ou \tech{app/globals.css}.

\section{Septième expertise}

\subsection{Contenu français}

\begin{description}[leftmargin=3.3cm,style=nextline]
  \item[Titre] Développement \& Prototypage
  \item[Sous-titre] Outils \& interfaces autour de la donnée
  \item[Sélection compacte] TypeScript · JavaScript · React · Next.js
\end{description}

La description précise que le domaine concerne les interfaces, outils internes, prototypes et
applications d'appui autour de la donnée. Elle indique explicitement que le portfolio démontre
le volet TypeScript, React et Next.js, sans transformer le titre professionnel en Software
Engineer ou Web Developer.

\subsection{Contenu anglais}

\begin{description}[leftmargin=3.3cm,style=nextline]
  \item[Title] Development \& Prototyping
  \item[Subtitle] Tools \& interfaces around data
  \item[Compact selection] TypeScript · JavaScript · React · Next.js
\end{description}

Le contenu anglais applique la même nuance de positionnement et les mêmes limites de
véracité que le contenu français.

\section{Organisation des technologies}

La carte détaillée organise les technologies en trois groupes conceptuels :

\begin{tabularx}{\textwidth}{>{\bfseries}p{0.29\textwidth}X}
  Programmation / Scripting & Python, JavaScript, TypeScript, VBA \\
  Web / UI & HTML5, CSS3, React, Next.js \\
  Outils de développement & Git, GitHub, VS Code \\
\end{tabularx}

La liste complète apparaît dans la description et sous forme de groupes de badges. Les autres
cartes conservent leur structure et leurs technologies existantes. Le modèle
\tech{ExpertiseItem} accepte désormais un sous-titre et des groupes d'outils optionnels ; les
six branches historiques continuent d'utiliser leur tableau plat sans migration inutile.

\section{Comportement du réseau desktop et tablette}

Le réseau SVG comporte désormais sept chemins réels depuis le nœud central
\textbf{DATA \& BUSINESS}. La nouvelle branche est placée dans la zone inférieure droite. La
branche Data Quality est légèrement remontée afin de conserver une respiration suffisante et
d'éviter les collisions de cartes ou de texte.

Chaque branche, y compris la nouvelle, conserve le contrat interactif existant :

\begin{itemize}
  \item tracé SVG neutre par défaut et doré lorsqu'il est actif ;
  \item animation initiale du tracé via \tech{pathLength} ;
  \item deux particules \tech{animateMotion} parcourant exactement le chemin associé ;
  \item activation synchronisée entre la carte détaillée, le chemin et le nœud ;
  \item activation au survol, au focus clavier et au clic ;
  \item attribut \tech{aria-pressed} reflétant l'état actif ;
  \item palette alimentée par les variables existantes en mode clair et sombre.
\end{itemize}

À 768 px, la visualisation SVG desktop/tablette reste affichée sous la liste des cartes. À partir
du breakpoint \tech{lg}, elle reprend sa place dans la deuxième colonne de la composition. Le
nœud central reste volontairement décalé très légèrement vers la gauche afin d'équilibrer les
trois branches gauches et les quatre branches droites.

\section{Comportement mobile}

Sous 768 px, l'écosystème simplifié conserve le centre, la spine verticale animée et les nœuds
en grille de deux colonnes. Comme le nombre de domaines est impair, le septième nœud
\emph{Développement \& Prototypage / Development \& Prototyping} occupe volontairement les
deux colonnes de la dernière ligne. Cette règle évite une cellule vide accidentelle et donne une
conclusion visuelle claire au réseau.

Les sept cartes détaillées restent empilées sur une colonne. Le mobile ne reçoit ni curseur
personnalisé ni tilt dépendant de la souris. Les boutons du réseau mobile répondent au focus, au
survol disponible et au clic, et exposent également leur état avec \tech{aria-pressed}.

\section{Identité visuelle, animation et accessibilité}

La version 0.7.1 conserve sans modification :

\begin{itemize}
  \item l'identité crème, navy et or ;
  \item les variables dédiées aux thèmes clair et sombre ;
  \item le mouvement du pattern, de l'aura et du ticker ;
  \item les transitions Motion, le tilt subtil et le curseur desktop ;
  \item le focus visible doré ;
  \item la suppression du curseur custom sur pointeur grossier ;
  \item la règle globale \tech{prefers-reduced-motion: reduce} ;
  \item les métriques Air France déjà documentées et vérifiées.
\end{itemize}

Aucune information essentielle ne dépend du mouvement : titres, descriptions, outils et état
actif restent compréhensibles lorsque les animations sont réduites.

\section{Règles de contenu et de positionnement}

Cette branche décrit une capacité complémentaire au travail Data et Business. Elle ne doit pas
devenir un changement de métier affiché. Les formulations futures doivent respecter les règles
suivantes :

\begin{enumerate}
  \item maintenir \textbf{Data \& Business Analyst} comme titre principal ;
  \item parler d'outils et d'interfaces \emph{autour de la donnée} ;
  \item ne pas attribuer un usage professionnel à React, Next.js ou TypeScript sans preuve ;
  \item pouvoir citer DAMERGI.COM comme démonstration directe de TypeScript, React et Next.js ;
  \item ne jamais inventer de mission, client, résultat, date ou métrique.
\end{enumerate}

\section{Matrice de vérification responsive}

\begin{tabularx}{\textwidth}{p{0.16\textwidth}p{0.25\textwidth}X}
  \textbf{Largeur} & \textbf{Mode} & \textbf{Contrôle Expertise 0.7.1} \\
  390 px & Smartphone étroit & Sept cartes lisibles ; dernier nœud mobile sur toute la ligne ; aucun débordement horizontal. \\
  430 px & Smartphone large & Groupes de badges lisibles ; rythme vertical et grille mobile équilibrés. \\
  768 px & Tablette portrait & Cartes sur une colonne ; réseau SVG complet et sept branches visibles sans chevauchement. \\
  Desktop & Deux colonnes & Nouvelle branche en bas à droite ; centre équilibré ; chemins et libellés séparés. \\
\end{tabularx}

Les vérifications doivent être répétées en français et en anglais, en clair et en sombre, au
clavier et avec la préférence de mouvement réduit.

\section{Fichiers concernés}

\begin{description}[leftmargin=5.4cm,style=nextline]
  \item[\tech{components/Expertise/ExpertiseSection.tsx}] septième clé, contenus FR/EN,
  groupes de technologies, chemin et position SVG, synchronisation interactive et traitement
  de la dernière ligne mobile.
  \item[\tech{docs/damergi\_documentation\_v0.7.1.tex}] présent document.
  \item[\tech{docs/damergi\_documentation\_v0.7.1.pdf}] version compilée lorsque LaTeX est
  disponible localement.
\end{description}

Le PDF historique \tech{docs/damergi\_documentation\_v0.7.0.pdf} est conservé intact.

\section{Historique des versions}

Les versions sont des jalons documentaires internes et ne supposent pas l'existence d'un tag
Git public.

\begin{tabularx}{\textwidth}{p{0.12\textwidth}p{0.21\textwidth}X}
  \textbf{Version} & \textbf{Étape} & \textbf{Principales modifications} \\
  0.1.0 & Bootstrap & Création Next.js et structure initiale. \\
  0.2.0 & Branding & Palette crème/navy/or, logo AFD, navbar, portrait, contacts et ticker. \\
  0.3.0 & Storytelling & Timeline narrative, portrait, pattern animé et aura. \\
  0.4.0 & Navigation & Timeline cliquable, FR/EN, thème, loader et favicon. \\
  0.5.0 & Détail & Modales, navigation et modes image/logo/tools. \\
  0.6.0 & Responsive & Hero, timeline et modales mobiles, CV et persistance. \\
  0.7.0 & Pré-déploiement & Expertise interactive à six branches, réseau animé, métriques, SEO, Open Graph, 404, robots et sitemap. \\
  \textbf{0.7.1} & \textbf{Expertise} & \textbf{Septième branche Développement \& Prototypage, réseau desktop rééquilibré, particules dédiées, technologies groupées, dernier nœud mobile pleine largeur et documentation du portfolio comme preuve TypeScript/React/Next.js.} \\
\end{tabularx}

\section{Continuité}

Les règles de maintenance de la version 0.7.0 restent en vigueur : modifier localement, tester
FR/EN, clair/sombre et les largeurs de référence, conserver les zones stables, exécuter les
contrôles TypeScript/lint/build, puis documenter chaque évolution significative. Le couple
\tech{damergi\_documentation\_v0.7.0.pdf} + présent document constitue la documentation de
référence de la version 0.7.1.

\vfill
\begin{center}
  {\color{damergigold}\rule{0.55\textwidth}{0.6pt}}\\[0.8em]
  \textbf{Fin de la documentation -- DAMERGI.COM -- Version \version}
\end{center}

\end{document}
